\section{讨论与局限性}
\label{sec:discussion}

\subsection{效率--精度权衡}
OmniSAM 的 memory mechanism 有助于连接全景视图，却使训练成本随历史 token 增长。MTC 与 SMR 分别压缩空间冗余和筛选内容相关记忆，以轻微 mIoU、BFScore 代价换取约 22.4\% 的训练时间下降。该结果说明，跨视场建模并不要求所有历史位置以原始密度参与每次 attention；结构化聚合和内容选择可以将更多计算集中在较紧凑的上下文上。消融结果还表明，降低 token budget 能减少 memory attention time，但边界和小目标上下文可能更敏感，且 total step time 不严格单调。原因在于完整训练还包含编码、解码、稀疏选择和数据处理，attention 降低不会按相同比例转化为单步时间收益。因此该方法应被理解为训练阶段的可控效率--精度调节，而非无损压缩；本文也不推断尚未验证的推理或显存收益。

\subsection{局限性}
实验集中于 Cityscapes $\rightarrow$ DensePASS、sam2\_s、13 类 single-seed 设置，尚不足以证明跨数据集或多随机种子稳定性。虽然 Baseline 与 Ours 使用相同协议，单次运行仍无法区分模块效应与随机波动，训练时间也可能受硬件状态影响。现有定性结果展示了四个全景样本的主要结构与局部差异，但样本数量有限，类别级、边界和小目标诊断仍不充分。本文只验证训练阶段收益，没有报告部署端推理速度或显存峰值。PQR 的证据同样有限：类别均衡过滤仅有轻微改善，原型门控未进一步提升，因而不能作为稳定性能来源。

\subsection{后续工作}
后续将优先补充多随机种子和一项额外全景数据集验证，以报告均值、方差和不同域偏移下的稳定性。细粒度分析将按类别、边界、ERP 接缝与小目标统计 token reduction 的影响，并结合更多 failure cases 判断被舍弃上下文的类型。在此基础上，可探索结合跨视场记忆一致性的伪标签可靠性建模，使过滤依据同时利用类别分布和历史视图；推理效率与显存占用则应在独立协议中评估，避免与训练收益混写。
